\Hefteintrag{1}{Was sind Cookies?}{%
\vspace{0.5cm}
\begin{itemize}
    \item Cookies sind kleine Datenelemente, die eine Website \emphColA{auf dem Computer}/im Browser des Nutzers \emphColA{speichert}.
    \item Beim Aufrufen einer Internetseite, schickt der Browser die gespeicherten Cookies an die Seite, sodass diese den Nutzer anhand der Cookies erkennen kann.
    \item Cookies sind nützlich, um bspw. \LoesungLuecke{Warenkörbe~}{6cm} und \LoesungLuecke{Logins~}{4cm} zu speichern. Diese würden ansonsten bei jedem Wechsel einer Unterseite verloren gehen.
    \item Mit Cookies können jedoch auch umfassende Nutzungsprofile erstellt werden, sodass Websites und Werbeanbieter genau herausfinden können, auf welchen Seiten sich ein Nutzer bewegt hat. 
    \item So erhält man personalisierte Werbung, kann jedoch auch \LoesungLuecke{gezielt manipuliert~}{6cm} werden.
    \item Werden Cookies im Browser gelöscht, stehen diese für Internetseiten nicht mehr zur Verfügung.
    \item Man unterscheidet zwischen: \begin{itemize}
        \item \LoesungLuecke{technisch notwendig}{14cm}
        \item \LoesungLuecke{intern für Komfort/Performance}{14cm}
        \item \LoesungLuecke{Drittanbieter (für Tracking)}{14cm}
    \end{itemize}
    \item Letztere sammeln hierbei Informationen für andere Internetseite (z.B. Werbe-Anbieter oder Online-Shops).
    \item Heutzutage gibt es viele Alternativen, um auch ohne Drittanbieter-Cookies umfassende Nutzerprofile zu erstellen. Dennoch sind Cookies hierzu weiterhin hoch-relevant.
\end{itemize}
}